% ============================================================
% EXAMEN PARCIAL — UNIDAD I · Paralelo B
% Ingeniería de Requisitos · UTEQ · 2025-2026
% ============================================================
\documentclass[11pt,a4paper]{article}

% ── Paquetes ─────────────────────────────────────────────────
\usepackage[T1]{fontenc}
\usepackage{babel}
\usepackage{geometry}
\geometry{left=2.2cm,right=2.2cm,top=1.8cm,bottom=2cm,headheight=22pt}
\usepackage{enumitem}
\usepackage{array,tabularx,booktabs,colortbl}
\usepackage{xcolor}
\usepackage{tikz}
\usepackage{fancyhdr}
\usepackage{graphicx}
\usepackage{setspace}
\usepackage{parskip}
\usepackage{multicol}
\usepackage{lastpage}

% ── Colores institucionales ─────────────────────────────────
\definecolor{uteqgreen}{HTML}{1B6B3A}
\definecolor{uteqgold}{HTML}{E8A317}
\definecolor{lightgray}{HTML}{F2F2F2}
\definecolor{darkgreen}{HTML}{1B5E30}
\definecolor{redborder}{HTML}{CC3333}

% ── Encabezado y pie ─────────────────────────────────────────
\pagestyle{fancy}
\fancyhf{}
\renewcommand{\headrulewidth}{2pt}
\renewcommand{\headrule}{\hbox to\headwidth{%
		\color{uteqgreen}\leaders\hrule height \headrulewidth\hfill}}
\lhead{\small\color{uteqgreen}\textbf{Ingeniería de Requisitos $\cdot$ Unidad I}}
\rhead{\small\color{uteqgreen}\textbf{UTEQ $\cdot$ 2025-2026}}
\cfoot{\thepage\ de \pageref{LastPage}}

% ── Comandos auxiliares ──────────────────────────────────────
\newcommand{\pts}[1]{%
	\hfill\tikz[baseline=(X.base)]{%
		\node[fill=uteqgold,text=white,font=\small\bfseries,
		rounded corners=3pt,inner sep=3pt](X){#1};}}

\newcommand{\lineas}[1]{%
	\par\vspace{4pt}%
	\foreach \i in {1,...,#1}{%
		.\,.\,.\,.\,.\,.\,.\,.\,.\,.\,.\,.\,.\,.\,.\,.\,.\,.\,%
		.\,.\,.\,.\,.\,.\,.\,.\,.\,.\,.\,.\,.\,.\,.\,.\,.\,.\,%
		.\,.\,.\,.\,.\,.\,.\,.\,.\,.\,.\,.\,.\,.\,.\,.\,.\,.\,%
		.\,.\,.\,.\,.\,.\,.\,.\,.\,.\,.\,.\,.\,.\,.\,.\,.\,.\,%
		.\,.\,.\,.\par\vspace{4pt}}}

\newcommand{\secbox}[1]{%
	\par\vspace{8pt}%
	\noindent\colorbox{uteqgreen}{%
		\parbox{\dimexpr\linewidth-2\fboxsep}{%
			\color{white}\bfseries\large #1}}%
	\par\vspace{6pt}}

% Macro para encabezado de tabla (evita Misplaced \noalign)
\newcommand{\thdr}[1]{\cellcolor{darkgreen}\color{white}\bfseries #1}

% ============================================================
\begin{document}
	
	% ── PORTADA ──────────────────────────────────────────────────
	\thispagestyle{fancy}
	\begin{center}
		\colorbox{uteqgreen}{\parbox{\linewidth}{\centering\color{white}\bfseries
				\vspace{6pt}
				UNIVERSIDAD TÉCNICA ESTATAL DE QUEVEDO\\[2pt]
				\small\itshape La primera universidad agropecuaria del Ecuador
				\vspace{6pt}}}
	\end{center}
	
	\vspace{1.5cm}
	\begin{center}
		{\LARGE\bfseries EXAMEN PARCIAL\\UNIDAD I}\\[10pt]
		{\large Fundamentos de Ingeniería de Requisitos\\y Proceso de Requisitos}
	\end{center}
	
	\vspace{1cm}
	\noindent
	\fbox{\parbox{\dimexpr\linewidth-2\fboxsep}{%
			\begin{tabularx}{\linewidth}{@{}>{\bfseries\color{uteqgreen}}l X@{}}
				Asignatura:        & Ingeniería de Requisitos \\[4pt]
				Resultado de       & \textit{Aplica los fundamentos teóricos de la IR} \\
				Aprendizaje:       & \textit{identificando tipos, propiedades y el proceso} \\
				& \textit{de requisitos para analizar su aplicabilidad} \\
				& \textit{en proyectos de software reales.} \\[4pt]
				Modalidad:         & Individual $\cdot$ Presencial \\
				Tiempo:            & 120 minutos \\
				Puntaje total:     & \textbf{10 puntos} \\
				Período:           & 2025--2026 \\
	\end{tabularx}}}
	
	\vspace{0.8cm}
	\noindent
	\fbox{\parbox{\dimexpr\linewidth-2\fboxsep}{%
			\begin{tabularx}{\linewidth}{@{}>{\bfseries}l X@{}}
				Apellidos y Nombres: & \hrulefill \\[8pt]
				Cédula / ID:         & \hrulefill \\[8pt]
				Paralelo:            & B \hspace{2cm} \textbf{Fecha:} \hrulefill \\
	\end{tabularx}}}
	
	% ── INSTRUCCIONES ────────────────────────────────────────────
	\newpage
	
	\noindent\fcolorbox{redborder}{white}{\parbox{\dimexpr\linewidth-2\fboxsep}{%
			{\color{redborder}\bfseries INSTRUCCIONES GENERALES --- Lea antes de comenzar:}
			\begin{enumerate}[leftmargin=*,itemsep=2pt]
				\item Lea cuidadosamente el caso de estudio completo antes de responder
				cualquier tarea.
				\item Todas las tareas se derivan del mismo contexto; sus respuestas deben
				ser coherentes entre sí.
				\item Fundamente cada respuesta con terminología técnica de la IR vista en
				la Unidad~I.
				\item Las respuestas sin justificación o que usen términos vagos sin
				cuantificar no tendrán puntaje.
				\item No se permite consulta de material externo, dispositivos electrónicos
				ni comunicación con terceros.
				\item Escriba con letra legible.\ Si usa tabla, respétela; si escribe
				párrafos, organícelos claramente.
	\end{enumerate}}}
	
	% ── CASO DE ESTUDIO ──────────────────────────────────────────
	\newpage
	
	\begin{center}
		{\large\bfseries CASO DE ESTUDIO:\\[4pt]
			SISTEMA DE GESTIÓN ACADÉMICA\\
			DEL INSTITUTO TECNOLÓGICO «SIMÓN BOLÍVAR»}
	\end{center}
	
	\vspace{6pt}
	\noindent\fbox{\parbox{\dimexpr\linewidth-2\fboxsep}{%
			\textbf{Contexto general}\\[4pt]
			El Instituto Tecnológico Superior «Simón Bolívar», ubicado en Babahoyo,
			Ecuador, atiende a \textbf{3\,200 estudiantes} distribuidos en 12 carreras
			tecnológicas. Actualmente, los procesos de matrícula, registro de
			calificaciones, control de asistencia y emisión de certificados se realizan
			de forma \textbf{manual mediante hojas de cálculo y documentos impresos},
			lo que genera duplicación de datos, errores en actas de notas, retrasos en
			la emisión de certificados y dificultad para generar reportes estadísticos
			para los organismos de control (SENESCYT y CES).
			
			La dirección del instituto ha decidido contratar el desarrollo de un
			\textbf{Sistema de Gestión Académica (SGA-SB)}. Usted ha sido designado
			como \textbf{analista líder de requisitos} del proyecto.
			
			Durante las primeras reuniones con los stakeholders, se recopilaron las
			siguientes declaraciones:}}
	
	\vspace{8pt}
	
	% ── Declaraciones de stakeholders ────────────────────────────
	\noindent\colorbox{uteqgreen!15}{\parbox{\dimexpr\linewidth-2\fboxsep}{%
			\textbf{\color{uteqgreen}Ing.\ Marlene Salazar --- Rectora del Instituto}}}
	
	\noindent\textit{«Necesito un sistema que me permita tomar decisiones con datos
		reales. Quiero ver en un panel cuántos estudiantes hay matriculados por
		carrera, la tasa de deserción y las calificaciones promedio por período.
		El sistema debe estar disponible siempre, especialmente en época de
		matrículas. No puede fallar.»}
	
	\vspace{6pt}
	\noindent\colorbox{uteqgreen!15}{\parbox{\dimexpr\linewidth-2\fboxsep}{%
			\textbf{\color{uteqgreen}Lic.\ Roberto Cedeño --- Coordinador Académico}}}
	
	\noindent\textit{«Los docentes deben poder registrar calificaciones parciales
		y finales con sus respectivos componentes ---prácticas, exámenes y
		proyectos--- y que el sistema calcule promedios automáticamente. Si un
		estudiante tiene una calificación reprobatoria, el sistema debe alertar
		al tutor asignado para que active el acompañamiento académico.»}
	
	\vspace{6pt}
	\noindent\colorbox{uteqgreen!15}{\parbox{\dimexpr\linewidth-2\fboxsep}{%
			\textbf{\color{uteqgreen}Prof.\ Diana Moreira --- Docente titular}}}
	
	\noindent\textit{«Yo necesito ver la lista de estudiantes de cada materia
		antes de la primera clase. Si un estudiante no está matriculado, no puedo
		aceptarlo. Y cuando registro notas, necesito que el sistema me confirme
		que se guardaron correctamente. Rápido, porque tengo muchas materias.»}
	
	\vspace{6pt}
	\noindent\colorbox{uteqgreen!15}{\parbox{\dimexpr\linewidth-2\fboxsep}{%
			\textbf{\color{uteqgreen}Ing.\ Carlos Zambrano --- Jefe de TI del Instituto}}}
	
	\noindent\textit{«El sistema debe desplegarse en la infraestructura local:
		tenemos un servidor Ubuntu 22.04~LTS con PostgreSQL~15. No podemos usar
		nube pública porque la normativa del CES para institutos tecnológicos
		exige que los datos académicos permanezcan en servidores nacionales.
		Además, debe integrarse con el sistema de cobros existente, el
		FinanCES~v1.8.»}
	
	\vspace{6pt}
	\noindent\colorbox{uteqgreen!15}{\parbox{\dimexpr\linewidth-2\fboxsep}{%
			\textbf{\color{uteqgreen}Ab.\ Patricia Loor --- Secretaria General}}}
	
	\noindent\textit{«El sistema debe cumplir con la normativa del CES y de
		SENESCYT para registro académico. Necesito generar actas de calificaciones
		con firma electrónica, certificados de matrícula y récords académicos
		que los estudiantes puedan descargar en PDF. Cada período reportamos
		indicadores al SIIES.»}
	
	\vspace{6pt}
	\noindent\colorbox{uteqgreen!15}{\parbox{\dimexpr\linewidth-2\fboxsep}{%
			\textbf{\color{uteqgreen}Sr.\ Kevin Astudillo --- Representante estudiantil}}}
	
	\noindent\textit{«Cuando me matriculo, quiero ver las materias disponibles
		según mi malla, saber los horarios y que no se crucen. Si ya aprobé una
		materia, que no me la ofrezca otra vez. Y quiero consultar mis notas
		desde el celular sin tener que ir a secretaría.»}
	
	% ── BORRADOR DE REQUISITOS ───────────────────────────────────
	\newpage
	
	\secbox{BORRADOR INICIAL DE REQUISITOS --- elaborado por un analista junior}
	
	\vspace{4pt}
	\noindent
	\begin{tabular}{@{}c p{13cm}@{}}
		\toprule
		\textbf{ID} & \textbf{Enunciado del requisito} \\
		\midrule
		R-01 & El sistema debe ser rápido al cargar las listas de estudiantes. \\[3pt]
		R-02 & El docente registrará las calificaciones parciales y finales de cada
		estudiante en la materia asignada. \\[3pt]
		R-03 & El sistema enviará una alerta al tutor cuando un estudiante obtenga
		una calificación reprobatoria. \\[3pt]
		R-04 & El sistema debe ser seguro y proteger los datos académicos. \\[3pt]
		R-05 & El sistema funcionará 24 horas al día, 7 días a la semana, sin
		interrupciones. \\[3pt]
		R-06 & El sistema debe ser compatible con Ubuntu 22.04~LTS y
		PostgreSQL~15. \\[3pt]
		R-07 & El sistema generará actas de calificaciones en formato PDF de forma
		automática. \\[3pt]
		R-08 & El sistema debe integrarse con el FinanCES~v1.8 para verificar pagos
		antes de permitir la matrícula. \\[3pt]
		R-09 & El estudiante podrá consultar sus calificaciones desde un dispositivo
		móvil. \\[3pt]
		R-10 & El sistema debe cumplir con la normativa del CES y SENESCYT para
		el registro académico de institutos tecnológicos. \\
		\bottomrule
	\end{tabular}
	
	% ============================================================
	%                    T A R E A S
	% ============================================================
	
	\newpage
	\secbox{TAREA --- Clasificación y Tipología del Sistema}
	
	% ── 1.1 ──────────────────────────────────────────────────────
	\subsection*{1.1~~Tipo de sistema \pts{0.5 pts}}
	
	\noindent\rule{4pt}{12pt}\hspace{4pt}%
	Clasifique el SGA-SB según la tipología de sistemas estudiada en la
	Unidad~I (\textit{sistema de información, embebido, software-intensivo o
		adaptativo}). Justifique su respuesta considerando las características
	del sistema y las declaraciones de los stakeholders.
	
	\lineas{4}
	
	% ── 1.2 ──────────────────────────────────────────────────────
	\subsection*{1.2~~Límite del sistema (\textit{system boundary}) \pts{0.5 pts}}
	
	\noindent\rule{4pt}{12pt}\hspace{4pt}%
	Identifique con precisión qué componentes, procesos y actores pertenecen
	al sistema y cuáles quedan fuera pero interactúan con él. Presente su
	respuesta en \textbf{dos listas diferenciadas}.
	
	\vspace{6pt}
	\noindent
	\begin{tabular}{|p{7cm}|p{7cm}|}
		\hline
		\thdr{Dentro del límite del sistema} &
		\thdr{Fuera del límite (interactúa con el sistema)} \\
		\hline
		& \\[4cm]
		\hline
	\end{tabular}
	
	% ── 1.3 ──────────────────────────────────────────────────────
	\subsection*{1.3~~Declaración de Visión \pts{0.5 pts}}
	
	\noindent\rule{4pt}{12pt}\hspace{4pt}%
	Redacte una declaración de \textbf{visión} del SGA-SB en no más de cinco
	líneas, que comunique: el objetivo estratégico, el valor que aportará y
	el problema que resuelve, conforme al concepto de visión del framework
	de IR.
	
	\lineas{5}
	
	% ── 1.4 ──────────────────────────────────────────────────────
	\newpage
	\subsection*{1.4~~Perspectivas de contexto \pts{0.5 pts}}
	
	\noindent\rule{4pt}{12pt}\hspace{4pt}%
	De las \textbf{ocho perspectivas de contexto} de la IR (uso, sujeto, TI,
	desarrollo, negocio, innovación, sociedad, técnica), identifique las
	\textbf{tres más relevantes} para este proyecto. Para cada una, indique
	un elemento concreto del caso que la activa y explique por qué es
	importante considerarla.
	
	\vspace{6pt}
	\noindent
	\begin{tabular}{|l|p{5cm}|p{5.5cm}|}
		\hline
		\thdr{Perspectiva} &
		\thdr{Elemento concreto del caso} &
		\thdr{Importancia para la IR} \\
		\hline
		& & \\[1.5cm]
		\hline
		& & \\[1.5cm]
		\hline
		& & \\[1.5cm]
		\hline
	\end{tabular}
	
	% ============================================================
	\newpage
	\secbox{TAREA --- Identificación y Clasificación de Stakeholders y Requisitos}
	
	% ── 2.1 ──────────────────────────────────────────────────────
	\subsection*{2.1~~Tabla de stakeholders \pts{0.5 pts}}
	
	\noindent\rule{4pt}{12pt}\hspace{4pt}%
	Elabore una tabla de stakeholders con las columnas indicadas. Identifique
	\textbf{al menos un conflicto real} entre dos stakeholders del caso.
	
	\vspace{6pt}
	\noindent
	{\small
		\begin{tabular}{|l|p{2.8cm}|l|p{3.8cm}|}
			\hline
			\thdr{Stakeholder} &
			\thdr{Rol en el proceso IR} &
			\thdr{Tipo de req.\ principal} &
			\thdr{Posible conflicto de interés} \\
			\hline
			& & & \\[0.9cm]\hline
			& & & \\[0.9cm]\hline
			& & & \\[0.9cm]\hline
			& & & \\[0.9cm]\hline
			& & & \\[0.9cm]\hline
			& & & \\[0.9cm]\hline
	\end{tabular}}
	
	% ── 2.2 ──────────────────────────────────────────────────────
	\subsection*{2.2~~Clasificación de requisitos extraídos \pts{0.5 pts}}
	
	\noindent\rule{4pt}{12pt}\hspace{4pt}%
	De las declaraciones de los stakeholders, extraiga y clasifique
	exactamente: \textbf{4~requisitos funcionales}, \textbf{3~requisitos no
		funcionales} y \textbf{2~restricciones del sistema}.
	
	\vspace{6pt}
	\noindent
	{\small
		\begin{tabular}{|c|p{5.5cm}|p{2.8cm}|p{2cm}|}
			\hline
			\thdr{ID} &
			\thdr{Enunciado (paráfrasis)} &
			\thdr{Stakeholder fuente} &
			\thdr{Clasificación} \\
			\hline
			RF-1  & & & \\[0.6cm]\hline
			RF-2  & & & \\[0.6cm]\hline
			RF-3  & & & \\[0.6cm]\hline
			RF-4  & & & \\[0.6cm]\hline
			NFR-1 & & & \\[0.6cm]\hline
			NFR-2 & & & \\[0.6cm]\hline
			NFR-3 & & & \\[0.6cm]\hline
			REST-1 & & & \\[0.6cm]\hline
			REST-2 & & & \\[0.6cm]\hline
	\end{tabular}}
	
	% ── 2.3 ──────────────────────────────────────────────────────
	\newpage
	\subsection*{2.3~~De necesidad de usuario a requisito cuantificable \pts{0.5 pts}}
	
	\noindent\rule{4pt}{12pt}\hspace{4pt}%
	La Prof.\ Diana Moreira indica: \textit{«Rápido, porque tengo muchas
		materias.»} Explique por qué es una \textbf{necesidad del usuario} y no
	un requisito bien formulado. Luego reescríbala como \textbf{requisito no
		funcional cuantificable} que cumpla verificabilidad y completitud.
	
	\vspace{6pt}
	\textbf{Explicación:}
	
	\lineas{3}
	
	\textbf{Requisito NFR reformulado:}
	
	\lineas{2}
	
	% ============================================================
	\newpage
	\secbox{TAREA --- Análisis Crítico del Borrador de Requisitos}
	
	% ── 3.1 ──────────────────────────────────────────────────────
	\subsection*{3.1~~Auditoría del borrador R-01 a R-10 \pts{1.5 pts}}
	
	\noindent\rule{4pt}{12pt}\hspace{4pt}%
	Analice cada requisito del borrador e identifique los problemas desde la
	perspectiva de las propiedades de un buen requisito. Para los que no
	presenten problemas, indíquelo explícitamente y justifique. Complete la
	tabla siguiente.
	
	\vspace{6pt}
	\noindent
	{\small
		\begin{tabular}{|c|p{3.2cm}|p{2.5cm}|p{5.5cm}|}
			\hline
			\thdr{ID} &
			\thdr{Problema detectado} &
			\thdr{Propiedad violada} &
			\thdr{Versión corregida y cuantificada} \\
			\hline
			R-01 & & & \\[0.7cm]\hline
			R-02 & & & \\[0.7cm]\hline
			R-03 & & & \\[0.7cm]\hline
			R-04 & & & \\[0.7cm]\hline
			R-05 & & & \\[0.7cm]\hline
			R-06 & & & \\[0.7cm]\hline
			R-07 & & & \\[0.7cm]\hline
			R-08 & & & \\[0.7cm]\hline
			R-09 & & & \\[0.7cm]\hline
			R-10 & & & \\[0.7cm]\hline
	\end{tabular}}
	
	% ── 3.2 ──────────────────────────────────────────────────────
	\newpage
	\subsection*{3.2~~Relación entre R-05 y R-06: restricción, dependencia o
		inconsistencia \pts{0.5 pts}}
	
	\noindent\rule{4pt}{12pt}\hspace{4pt}%
	Analice el par R-05 («funcionará 24/7 sin interrupciones») y R-06
	(«compatible con Ubuntu 22.04~LTS»). ¿Qué propiedad del conjunto de
	requisitos podría verse comprometida si el servidor requiere reinicios
	para actualizaciones de seguridad del kernel? Explique la relación entre
	ambos e indique si existe una dependencia, una restricción o una
	inconsistencia potencial.
	
	\lineas{4}
	
	% ── 3.3 ──────────────────────────────────────────────────────
	\subsection*{3.3~~Tratamiento correcto de R-10 en el proceso de IR \pts{0.5 pts}}
	
	\noindent\rule{4pt}{12pt}\hspace{4pt}%
	El requisito R-10 es técnicamente real pero, como está redactado, presenta
	un problema fundamental. Identifique el problema, clasifíquelo
	correctamente (funcional, no funcional o restricción) y proponga cómo
	debería tratarse en el proceso de IR para que sea accionable para el
	equipo de desarrollo.
	
	\lineas{6}
	
	% ============================================================
	\newpage
	\secbox{TAREA --- Selección y Justificación del Proceso de Requisitos}
	
	% ── 4.1 ──────────────────────────────────────────────────────
	\subsection*{4.1~~Cuadro comparativo de modelos del proceso de requisitos
		\pts{1.0 pts}}
	
	\noindent\rule{4pt}{12pt}\hspace{4pt}%
	Compare los tres modelos (cascada, iterativo, ágil) evaluando los
	criterios indicados \textbf{en el contexto específico del SGA-SB}.
	Concluya con una recomendación fundamentada del modelo más adecuado,
	citando al menos \textbf{tres razones} basadas en el caso.
	
	\vspace{6pt}
	\noindent
	\begin{tabular}{|p{3.5cm}|p{3.5cm}|p{3.5cm}|p{3.5cm}|}
		\hline
		\thdr{Criterio de evaluación} &
		\thdr{Cascada} &
		\thdr{Iterativo} &
		\thdr{Ágil} \\
		\hline
		Gestión de cambios de requisitos & & & \\[1cm]\hline
		Participación de los stakeholders del instituto & & & \\[1cm]\hline
		Adecuación al tipo de sistema (SGA-SB) & & & \\[1cm]\hline
		Riesgo principal en este contexto & & & \\[1cm]\hline
		Nivel de documentación requerido & & & \\[1cm]\hline
	\end{tabular}
	
	\vspace{6pt}
	\textbf{Recomendación fundamentada (modelo seleccionado y tres razones):}
	
	\lineas{5}
	
	% ── 4.2 ──────────────────────────────────────────────────────
	\newpage
	\subsection*{4.2~~Gestión de un cambio de alcance emergente \pts{0.5 pts}}
	
	\noindent\rule{4pt}{12pt}\hspace{4pt}%
	El Ing.\ Zambrano señala a mitad de la reunión: \textit{«Ah, y también
		necesitamos que el sistema pueda exportar datos en formato XML para el
		reporte semestral al SIIES de SENESCYT, porque eso nos piden cada
		semestre.»} Desde el \textbf{framework de IR}, identifique:
	\begin{enumerate}[label=(\alph*),itemsep=2pt]
		\item en qué actividad del proceso se detectó este nuevo requisito,
		\item a qué tipo de artefacto debe incorporarse, y
		\item qué actividad transversal del framework debe activarse.
	\end{enumerate}
	
	\lineas{3}
	
	% ── 4.3 ──────────────────────────────────────────────────────
	\subsection*{4.3~~Propiedades emergentes del SGA-SB \pts{0.5 pts}}
	
	\noindent\rule{4pt}{12pt}\hspace{4pt}%
	Identifique \textbf{dos propiedades emergentes} del SGA-SB. Para cada
	una: nómbrela, explique por qué es emergente y describa qué combinación
	de componentes la genera.
	
	\vspace{6pt}
	\textbf{1) Nombre:} \dotfill
	
	\textbf{Por qué es emergente:} \dotfill
	
	\hphantom{\textbf{Por qué es emergente:}} \dotfill
	
	\textbf{Componentes que la generan:} \dotfill
	
	\vspace{8pt}
	\textbf{2) Nombre:} \dotfill
	
	\textbf{Por qué es emergente:} \dotfill
	
	\hphantom{\textbf{Por qué es emergente:}} \dotfill
	
	\textbf{Componentes que la generan:} \dotfill
	
	% ============================================================
	\newpage
	\secbox{TAREA --- Aplicación de ISO/IEC 25010 y Síntesis}
	
	% ── 5.1 ──────────────────────────────────────────────────────
	\subsection*{5.1~~Especificación de NFR según ISO/IEC 25010 \pts{1.0 pts}}
	
	\noindent\rule{4pt}{12pt}\hspace{4pt}%
	Especifique cuatro requisitos no funcionales del SGA-SB, uno por cada
	característica de calidad indicada. Cada requisito debe ser
	\textbf{cuantificable y verificable}.
	
	\vspace{6pt}
	\noindent
	{\small
		\begin{tabular}{|p{2.2cm}|p{2cm}|p{4.5cm}|p{4.5cm}|}
			\hline
			\thdr{\scriptsize Característica ISO 25010} &
			\thdr{Subcategoría} &
			\thdr{Requisito NFR cuantificado} &
			\thdr{Criterio de verificación} \\
			\hline
			Fiabilidad & & & \\[1.2cm]\hline
			Seguridad & & & \\[1.2cm]\hline
			Usabilidad & & & \\[1.2cm]\hline
			Eficiencia de desempeño & & & \\[1.2cm]\hline
	\end{tabular}}
	
	% ── 5.2 ──────────────────────────────────────────────────────
	\subsection*{5.2~~Resolución de conflicto de prioridades entre stakeholders
		\pts{0.5 pts}}
	
	\noindent\rule{4pt}{12pt}\hspace{4pt}%
	La Ing.\ Salazar enfatiza la disponibilidad permanente y el panel de
	indicadores en tiempo real; la Prof.\ Moreira enfatiza la velocidad del
	registro de notas y la confirmación inmediata. Desde la perspectiva del
	proceso de requisitos y el rol del analista, explique qué debe hacer el
	analista ante este conflicto y qué actividad del proceso de IR se activa
	para resolverlo.
	
	\lineas{5}
	
	% ── 5.3 ──────────────────────────────────────────────────────
	\subsection*{5.3~~Reflexión y argumento técnico final \pts{0.5 pts}}
	
	\noindent\rule{4pt}{12pt}\hspace{4pt}%
	Un colega afirma: \textit{«Para un instituto tecnológico pequeño como
		este, no tiene sentido hacer todo este proceso de IR. Con preguntar a la
		rectora qué quiere y empezar a programar es suficiente.»} Elabore un
	argumento técnico fundamentado (\textbf{mínimo 150 palabras}) que refute
	esta afirmación, considerando: el tipo de sistema, las consecuencias de
	requisitos deficientes en sistemas académicos críticos, la calidad del
	proceso de IR y el impacto sobre la calidad del producto final.
	
	\lineas{9}
	
	% ============================================================
	\newpage
	\secbox{RÚBRICA DE EVALUACIÓN}
	
	\vspace{6pt}
	\noindent
	{\small
		\begin{tabular}{|c|p{3.8cm}|p{5.5cm}|c|}
			\hline
			\thdr{Tarea} &
			\thdr{Capacidad evaluada} &
			\thdr{Indicadores de logro} &
			\thdr{Puntaje} \\
			\hline
			T1 & Clasifica el sistema, define el límite, formula la visión e identifica
			perspectivas de contexto &
			Tipo de sistema correcto y justificado; límite preciso y diferenciado;
			visión alineada al RA; perspectivas pertinentes con evidencia del caso &
			2.0 \\
			\hline
			T2 & Identifica stakeholders, clasifica requisitos y transforma
			necesidades en NFR cuantificables &
			Tabla de stakeholders completa con conflictos identificados;
			clasificación RF/NFR/restricción correcta; NFR con métricas verificables &
			1.5 \\
			\hline
			T3 & Detecta deficiencias en requisitos, aplica propiedades de calidad y
			propone correcciones &
			Propiedad violada identificada con precisión; corrección cuantificada;
			análisis de relación R-05/R-06; tratamiento adecuado de R-10 &
			2.5 \\
			\hline
			T4 & Selecciona y justifica el modelo de proceso; aplica el framework de
			IR ante cambios &
			Cuadro comparativo con criterios contextualizados al SGA-SB;
			recomendación con 3 razones válidas; actividades del framework
			correctamente identificadas; propiedades emergentes justificadas &
			2.0 \\
			\hline
			T5 & Especifica NFR según ISO/IEC 25010; resuelve conflictos entre
			stakeholders; argumenta el valor de la IR &
			NFR cuantificados por característica con criterio de verificación;
			actividad de resolución correcta; argumento técnico coherente
			$\geq$~150 palabras &
			2.0 \\
			\hline
			\multicolumn{3}{|r|}{\cellcolor{darkgreen}\color{white}\bfseries TOTAL} &
			\cellcolor{darkgreen}\color{white}\bfseries 10.0 \\
			\hline
	\end{tabular}}
	
	\vspace{10pt}
	\textbf{Criterios transversales de calidad --- aplican a todas las tareas:}
	
	\begin{itemize}[leftmargin=*,itemsep=2pt]
		\item \textbf{Pertinencia:} Las respuestas deben usar el contexto del
		SGA-SB, no ejemplos genéricos.
		\item \textbf{Fundamentación:} Uso correcto y preciso de la terminología
		técnica de la IR (Unidad~I).
		\item \textbf{Coherencia:} Las respuestas de distintas tareas deben ser
		consistentes entre sí.
		\item \textbf{Cuantificación:} Todo requisito formulado o reformulado debe
		incluir métricas verificables. Términos vagos como «rápido»,
		«seguro» o «confiable» sin cuantificar reducen el puntaje
		proporcionalmente.
		\item \textbf{Extensión:} Ni tan breve que carezca de sustento, ni tan
		extensa que exceda el tiempo disponible.
	\end{itemize}
	
	% ── PIE FINAL ────────────────────────────────────────────────
	\vfill
	\begin{center}
		\rule{5cm}{0.4pt}\hspace{3cm}\rule{5cm}{0.4pt}\\[4pt]
		{\small Firma del Docente \hspace{3.5cm} Revisado por}\\
		{\small Ingeniería de Requisitos \hspace{2cm} Coordinación de Carrera}
	\end{center}
	
	\vspace{6pt}
	\begin{center}
		{\footnotesize Universidad Técnica Estatal de Quevedo $\cdot$ Facultad de
			Ciencias de la Ingeniería $\cdot$ Carrera de Ingeniería de Software
			$\cdot$ Período 2025--2026}
	\end{center}
	
\end{document}